\chapter{地磁気}
\section{磁場の方向と地理的な北}
\label{theory:geomagnetism:declination}
方位磁針が向く方向は地理的な真北と一般に一致しない。地磁気の水平成分を真北方向\(X\)、東方向\(Y\)に分けると、東偏を正とする規約では偏角\(D\)を\(D=\operatorname{atan2}(Y,X)\)で定める。国土地理院の磁気図で西偏を正の大きさとして表示する場合は符号規約が異なる。角の大きさ、東西方向、符号を別々に確かめる。

地球中心の理想軸双極子だけなら非常に単純な分布となるが、実際の磁場には非双極子成分と局所的な磁気異常がある。同じ緯度又は同じ経度にあるだけでは偏角は同一にならない。また地球内部の流体運動などに伴って永年変化するため、磁気図には基準時刻である元期が付される。

令和元年度問27の「現在」は2019年の試験時点である。試験前に公表された国土地理院の2015.0年磁気図を参照すると、日本周辺は西偏で、沖縄本島付近は約5度、北海道では概ね9〜10度である\cite{gsi_declination2015}。これらは当時の問題を判断するための概略値であり、2026年現在の各地点の測量値を表すものではない。緯度だけの普遍的単調則を主張せず、地図上の分布に基づいて誤答肢を判断する。

\section{誘導磁化・残留磁化と磁気探査}
\label{theory:geomagnetism:magnetization}
磁気探査は岩石や埋設物の磁化が生む磁場の差を測り、地下構造を推定する。岩石の磁化には、現在の地磁場に応答する誘導磁化と、過去の場などを記録する残留磁化がある。線形近似では
\[
 \boldsymbol{M}=\kappa\boldsymbol{H}+\boldsymbol{M}_r
\]
である。\(\kappa\)は磁化率、\(\boldsymbol{H}\)は磁場の強さ、\(\boldsymbol{M}_r\)は自然残留磁化である。両者の方向が一致するとは限らないので、磁気異常の正負を磁性物質の有無だけで決めることはできない。

ケーニヒスベルガー比は両者の大きさを比べる無次元量であり、
\[
 Q=\frac{|\boldsymbol{M}_r|}{|\kappa\boldsymbol{H}|}
\]
と書ける。\(Q\ll1\)なら誘導磁化が支配的、\(Q\gg1\)なら残留磁化を無視できない。地質調査研究報告の花崗岩の物理定数研究にもこの定義が使われている\cite{kanaya_okuma2011}。比を知っても磁化方向まで確定するわけではなく、密度の比、磁化率そのもの、地震波速度の比ではない。

強磁性・フェリ磁性体の自発磁化は熱運動による磁気秩序の変化と結び付いている。キュリー温度は自発磁化が消失する転移温度であり、磁性鉱物や磁気を帯びる層の解釈に関係する。一方、プランク関数は熱平衡放射のスペクトルを表す。岩石の温度を扱うという広い共通点だけで、磁気探査の基本用語と同列にしてはならない。

\section{探査データから製品名まで断定できるか}
磁化が同じでも形状・深度が違えば類似した磁気異常を生みうる。位置の概略と磁性体の存在を探ることと、細かな製品種類まで一意に識別することの間には逆問題の非一意性がある。平成23年度問29はこの通常の磁気探査上の限界を述べる。他の観測や形状モデルを組み合わせた現代の識別技術全般まで否定するものと読まない。共通理論は物理探査章でも観測量の違いから整理する。
\subsection{令和元年度 問27}
\label{q:r01-reg-applied-27}

\noindent\textbf{原典：} sources/R01/17\_applied\_science.pdf、PDF物理9頁。

\textbf{公式正答：} 2（single）。

\par\textit{原典本文の校正と独自解説は未作成。分類のみの登録項目。}

